\documentclass[11pt,a4paper]{article}
\usepackage[utf8]{inputenc}
\usepackage[T1]{fontenc}
\usepackage[spanish,es-nodecimaldot]{babel}
\usepackage{lmodern,microtype}
\usepackage[a4paper,landscape,margin=1.2cm]{geometry}
\usepackage{longtable}[=v4.13]
\usepackage{booktabs,array,xcolor,enumitem}
\usepackage[unicode,colorlinks=true,urlcolor=blue!45!black]{hyperref}
\pagestyle{empty}
\setlength{\parindent}{0pt}
\setlength{\tabcolsep}{4pt}
\setlength{\LTpre}{0pt}
\setlength{\LTpost}{0pt}
\setlength{\emergencystretch}{2em}
\renewcommand{\arraystretch}{1.04}
\newcolumntype{L}[1]{>{\raggedright\arraybackslash}p{#1}}
\makeatletter
\newcommand{\sourceurl}[2]{\hyper@linkurl{#2}{#1}}
\makeatother
\begingroup\catcode`\%=12
\expandafter\gdef\csname srcS01\endcsname{\sourceurl{file:///C:/Users/orlan/Desktop/SanBartolo/global_economics/book2026-2027/seminars/term_1/second_learning_evidence/C5%2C6%2C7%2C8%2C9%2C10.%20Nicol%C3%A1s%20%C3%81lvarez%2C%20Thomas%20%C3%81lvarez%2C%20Juan%20Felipe%20Ojeda..pdf}{Álvarez y Ojeda, C5--C10}}
\endgroup
\begingroup\catcode`\%=12
\expandafter\gdef\csname srcS02\endcsname{\sourceurl{file:///C:/Users/orlan/Desktop/SanBartolo/global_economics/book2026-2027/seminars/term_1/second_learning_evidence/Criterio_10.pdf}{Ayala y Galán, acumulativo C3--C10}}
\endgroup
\begingroup\catcode`\%=12
\expandafter\gdef\csname srcS03\endcsname{\sourceurl{file:///C:/Users/orlan/Desktop/SanBartolo/global_economics/book2026-2027/seminars/term_1/second_learning_evidence/Encuesta_SBLM%20%282%29.pdf}{Ayala y Galán, Encuesta SBLM}}
\endgroup
\begingroup\catcode`\%=12
\expandafter\gdef\csname srcS04\endcsname{\sourceurl{file:///C:/Users/orlan/Desktop/SanBartolo/global_economics/book2026-2027/seminars/term_1/second_learning_evidence/c5c6c7c8.pdf}{González y Barón, C5--C6}}
\endgroup
\begingroup\catcode`\%=12
\expandafter\gdef\csname srcS05\endcsname{\sourceurl{file:///C:/Users/orlan/Desktop/SanBartolo/global_economics/book2026-2027/seminars/term_1/second_learning_evidence/Proyecto_C7-C8.pdf}{González y Barón, Proyecto C7--C8}}
\endgroup
\begingroup\catcode`\%=12
\expandafter\gdef\csname srcS06\endcsname{\sourceurl{file:///C:/Users/orlan/Desktop/SanBartolo/global_economics/book2026-2027/seminars/term_1/second_learning_evidence/Lucio_IsaC_Criterio_5-10.pdf}{Isa y Lucio, hábitos de sueño}}
\endgroup
\begingroup\catcode`\%=12
\expandafter\gdef\csname srcS07\endcsname{\sourceurl{file:///C:/Users/orlan/Desktop/SanBartolo/global_economics/book2026-2027/seminars/term_1/second_learning_evidence/SeminariosMiguel_Ignacio_tex%20%283%29.pdf}{Caballero y Rodríguez, proyecto}}
\endgroup
\begingroup\catcode`\%=12
\expandafter\gdef\csname srcS08\endcsname{\sourceurl{file:///C:/Users/orlan/Desktop/SanBartolo/global_economics/book2026-2027/seminars/term_1/second_learning_evidence/C1-C10%2C%20Sebastian%20Lozano%2C%20Alejandro%20Camargo.pdf}{Lozano y Camargo, C1--C10}}
\endgroup
\begingroup\catcode`\%=12
\expandafter\gdef\csname srcS09\endcsname{\sourceurl{file:///C:/Users/orlan/Desktop/SanBartolo/global_economics/book2026-2027/seminars/term_1/second_learning_evidence/Segunda_evidencia_de_apredizaje__C6_C10_.pdf}{Del Cairo y Clavijo, residuos}}
\endgroup
\begingroup\catcode`\%=12
\expandafter\gdef\csname srcS10\endcsname{\sourceurl{file:///C:/Users/orlan/Desktop/SanBartolo/global_economics/book2026-2027/seminars/term_1/second_learning_evidence/Di%CC%81az%2C%20Maldonado%20%26%20Lo%CC%81pez%20C5%20%26%20C6.pdf}{Díaz, Maldonado y López, C5--C6}}
\endgroup
\begingroup\catcode`\%=12
\expandafter\gdef\csname srcS11\endcsname{\sourceurl{file:///C:/Users/orlan/Desktop/SanBartolo/global_economics/book2026-2027/seminars/term_1/second_learning_evidence/C3%2C%20Andre%CC%81s%20Maldonado%2C%20Carlos%20Lo%CC%81pez%2C%20Thomas%20Di%CC%81az.pdf}{Encuesta de graderías, 30 estudiantes}}
\endgroup
\begingroup\catcode`\%=12
\expandafter\gdef\csname srcS12\endcsname{\sourceurl{file:///C:/Users/orlan/Desktop/SanBartolo/global_economics/book2026-2027/seminars/term_1/second_learning_evidence/Lo%CC%81pez%2C%20Di%CC%81az%20%26%20Maldonado%20C7%20%26%20C8.docx}{López, Díaz y Maldonado, C7--C8 (DOCX)}}
\endgroup
\begingroup\catcode`\%=12
\expandafter\gdef\csname srcS13\endcsname{\sourceurl{file:///C:/Users/orlan/Desktop/SanBartolo/global_economics/book2026-2027/seminars/term_1/second_learning_evidence/Maldonado%2C%20Lopez%2C%20Diaz%20C9%20%26%20C10%20.docx}{Maldonado, López y Díaz, C9--C10 (DOCX)}}
\endgroup
\begingroup\catcode`\%=12
\expandafter\gdef\csname srcS14\endcsname{\sourceurl{file:///C:/Users/orlan/Desktop/SanBartolo/global_economics/book2026-2027/seminars/term_1/second_learning_evidence/Seminarios%20%281%29.pdf}{Guampe y Ospina, cafetería}}
\endgroup
\begingroup\catcode`\%=12
\expandafter\gdef\csname srcS15\endcsname{\sourceurl{file:///C:/Users/orlan/Desktop/SanBartolo/global_economics/book2026-2027/seminars/term_1/second_learning_evidence/Seminarios_2026_2027%20%282%29.pdf}{Julián Maldonado, organización académica}}
\endgroup
\begingroup\catcode`\%=12
\expandafter\gdef\csname srcS16\endcsname{\sourceurl{file:///C:/Users/orlan/Desktop/SanBartolo/global_economics/book2026-2027/seminars/term_1/second_learning_evidence/C5.pdf}{Ospina y Méndez, C5--C6}}
\endgroup
\begingroup\catcode`\%=12
\expandafter\gdef\csname srcS17\endcsname{\sourceurl{file:///C:/Users/orlan/Desktop/SanBartolo/global_economics/book2026-2027/seminars/term_1/second_learning_evidence/C7..pdf}{Ospina y Méndez, C7--C8}}
\endgroup
\begingroup\catcode`\%=12
\expandafter\gdef\csname srcS18\endcsname{\sourceurl{file:///C:/Users/orlan/Desktop/SanBartolo/global_economics/book2026-2027/seminars/term_1/second_learning_evidence/C9.pdf}{Ospina y Méndez, C9--C10}}
\endgroup
\begingroup\catcode`\%=12
\expandafter\gdef\csname srcS19\endcsname{\sourceurl{file:///C:/Users/orlan/Desktop/SanBartolo/global_economics/book2026-2027/seminars/term_1/second_learning_evidence/pedro_y_tape.pdf}{Pedro y Santiago, baños}}
\endgroup
\begingroup\catcode`\%=12
\expandafter\gdef\csname srcS20\endcsname{\sourceurl{file:///C:/Users/orlan/Desktop/SanBartolo/global_economics/book2026-2027/seminars/term_1/second_learning_evidence/Seminarios%20Tom%C3%A1s%20Rinc%C3%B3n.pdf}{Tomás Rincón, filas}}
\endgroup
\makeatletter
\newcommand{\srclink}[1]{\begingroup\renewcommand{\sourceurl}[2]{\hyper@linkurl{#1}{##1}}\csname src#1\endcsname\endgroup}
\makeatother

% Column fractions apply to the usable width after the eight internal paddings.
\newlength{\usablewidth}
\setlength{\usablewidth}{\dimexpr\textwidth-8\tabcolsep\relax}
\newenvironment{compactitems}{\begin{minipage}[t]{\linewidth}\raggedright\begin{itemize}[label=\textbullet,leftmargin=1em,labelsep=.35em,nosep,topsep=0pt,partopsep=0pt,parsep=0pt]}{\end{itemize}\end{minipage}}
\newcommand{\evid}[2]{(\srclink{#1}, pp.~#2)}
\begin{document}
\begin{longtable}{@{}L{.05\usablewidth}L{.125\usablewidth}L{.50\usablewidth}L{.13\usablewidth}L{.195\usablewidth}@{}}
\toprule
\textbf{Grupo} & \textbf{Integrantes} & \textbf{Trabajo C5--C10} & \textbf{Problema} & \textbf{Soluciones} \\
\midrule
\endfirsthead
\toprule
\textbf{Grupo} & \textbf{Integrantes} & \textbf{Trabajo C5--C10} & \textbf{Problema} & \textbf{Soluciones} \\
\midrule
\endhead
\bottomrule
\endfoot
\bottomrule
\endlastfoot
1 & \begin{compactitems}
\item Thomas Sebastián Álvarez Rey
\item Nicolás Álvarez Ospina
\item Juan Felipe Ojeda Hurtado
\end{compactitems} & \begin{compactitems}
\item \textbf{C5:} Organizaron el recorrido de compra y uso de microondas y tres hallazgos sobre colados, tiempo perdido y falta de regulación. \evid{S01}{1--2}
\item \textbf{C6:} Compararon filas, presión en debates y extensión de guías, y priorizaron las filas con 9 puntos frente a 7 y 6. \evid{S01}{2--4}
\item \textbf{C7:} Reformularon el desafío para garantizar turnos equitativos y reducir la espera al comprar alimentos o calentar loncheras. \evid{S01}{4--5}
\item \textbf{C8:} Relacionaron normas, acción colectiva y señalética; quedó pendiente verificar el peso de la confusión espacial. \evid{S01}{5}
\item \textbf{C9:} Registraron once ideas, incluidas una fila delimitada con supervisión y una zona de microondas con fichas. \evid{S01}{6--8}
\item \textbf{C10:} Contrastaron no atender a colados y supervisar la fila, con supuestos pendientes sobre autorización, aceptación de monitores y protección del descanso. \evid{S01}{8}
\end{compactitems} & Colados y desorden al comprar alimentos o usar microondas reducen el descanso. & \begin{compactitems}
\item Proponen delimitar la fila con supervisión estudiantil y no atender a colados, como posibilidad por explorar. \evid{S01}{7--8}
\item Proponen una zona independiente de microondas con fichas de tiempo por usuario. \evid{S01}{7}
\end{compactitems} \\
\midrule
2 & \begin{compactitems}
\item Patrik Ayala Palenik
\item Gabriel Hernando Galán Osorio
\end{compactitems} & \begin{compactitems}
\item \textbf{C5:} Reunieron perfiles, esperas de 5 a 9 minutos y observaciones y testimonios sobre filas, colados y conflictos. \evid{S02}{1--2}
\item \textbf{C6:} Identificaron reducir la espera como necesidad principal; no se identificó una matriz comparativa de oportunidades en los PDF revisados. \evid{S02}{2}; \evid{S03}{1--3}
\item \textbf{C7:} Revisaron el desafío para reducir la espera y los conflictos por colados de décimo al comprar o reclamar durante el descanso. \evid{S02}{2--3}
\item \textbf{C8:} Conectaron diseño de servicios y flujos físicos con convivencia y resolución pacífica de conflictos. \evid{S02}{3}
\item \textbf{C9:} Documentaron cinco propuestas, preguntas y ampliaciones sobre fila exprés, señalización, campaña y combinaciones. \evid{S02}{4}
\item \textbf{C10:} Contrastaron fila exprés y turnos por colores, examinando pago anticipado, respeto de marcas, sencillez y acceso sin teléfono. \evid{S02}{4--5}
\end{compactitems} & Espera y conflictos por colados al comprar o reclamar en Asobartolina. & \begin{compactitems}
\item Proponen fila exprés para reclamar pedidos prepagados. \evid{S02}{4}
\item Proponen turnos por colores, con sencillez y acceso sin teléfono como condiciones. \evid{S02}{4--5}
\item Proponen campaña gráfica de respeto de la fila. \evid{S02}{4}
\end{compactitems}
{\footnotesize Fuente complementaria: \srclink{S03}.} \\
\midrule
3 & \begin{compactitems}
\item Sebastián González Mojica
\item Emilio José Barón Romero
\end{compactitems} & \begin{compactitems}
\item \textbf{C5:} Organizaron el recorrido del uso de audífonos y cuatro hallazgos sobre ocultamiento, trabajo individual y beneficios percibidos diversos. \evid{S04}{3--4}
\item \textbf{C6:} Compararon claridad normativa, dinamismo pedagógico y acuerdos tecnológicos, y priorizaron aclarar el uso de audífonos con 9 puntos. \evid{S04}{5--8}
\item \textbf{C7:} Reformularon el desafío sobre gestionar la concentración en trabajo autónomo manteniendo la comunicación docente. \evid{S05}{2--3}
\item \textbf{C8:} Vincularon autorregulación y acuerdos de convivencia, distinguiendo conceptos de hipótesis y verificaciones institucionales pendientes. \evid{S05}{4--5}
\item \textbf{C9:} No se identificó un registro de ideación en los dos PDF accesibles y legibles revisados. \evid{S04}{1--8}; \evid{S05}{1--6}
\item \textbf{C10:} No se identificó contraste de ideas con evidencia del usuario en los dos PDF revisados. \evid{S04}{1--8}; \evid{S05}{1--6}
\end{compactitems} & Gestionar la concentración en trabajo autónomo, considerando audífonos, normas y comunicación docente; beneficios de la música no universales. & \begin{compactitems}
\item No se identificaron soluciones desarrolladas en los PDF revisados; autorregulación y acuerdos en C8 son conocimientos de apoyo. \evid{S04}{1--8}; \evid{S05}{2--6}
\end{compactitems} \\
\midrule
4 & \begin{compactitems}
\item Ana Isabella Colorado Parrado
\item Lucio Barrero Sánchez
\end{compactitems} & \begin{compactitems}
\item \textbf{C5:} Organizaron diez respuestas, el recorrido del estudiante y hallazgos sobre bienestar y horarios, con una tabla parcialmente legible. \evid{S06}{1--4}
\item \textbf{C6:} Compararon calidad del sueño y horario estable, y eligieron el horario por 8 a 6, con puntajes parcialmente justificados. \evid{S06}{4--5}
\item \textbf{C7:} Formularon y revisaron un desafío sobre hábitos de sueño de estudiantes de décimo con mala calidad de descanso. \evid{S06}{6}
\item \textbf{C8:} Relacionaron administración del tiempo y psicología de hábitos; quedaron pendientes los beneficios de mantener horarios estables. \evid{S06}{6--7}
\item \textbf{C9:} Registraron seis ideas y las combinaron en un horario sin pantallas y un planeador con relajación; algunas tablas son parcialmente legibles. \evid{S06}{7--10}
\item \textbf{C10:} Examinaron ambas propuestas mediante supuestos, preguntas pendientes y adaptación a responsabilidades y preferencias personales. \evid{S06}{8--9}
\end{compactitems} & Horarios de sueño irregulares, tareas y pantallas se relacionan con cansancio escolar. & \begin{compactitems}
\item Proponen calcular tiempos de tareas y organizar un horario personal con apagado de pantallas antes de dormir. \evid{S06}{7}
\item Proponen un planeador de obligaciones con relajación al terminar las tareas. \evid{S06}{8}
\end{compactitems} \\
\midrule
5 & \begin{compactitems}
\item Miguel Ángel Caballero Hernández
\item Juan Ignacio Rodríguez Rodríguez
\end{compactitems} & \begin{compactitems}
\item \textbf{C5:} Organizaron el recorrido de búsqueda y cuatro hallazgos con quince encuestados y registros de observación. \evid{S07}{3--5}
\item \textbf{C6:} Eligieron clasificación física frente a reducción de hurto o pérdida por 9 a 7, reconociendo que no resolvería desapariciones. \evid{S07}{5--6}
\item \textbf{C7:} Ajustaron el desafío para facilitar la búsqueda visual y física de pertenencias de Ciclo V sin desordenar la mesa. \evid{S07}{6--7}
\item \textbf{C8:} Relacionaron significantes visuales y metodología 5S para comprender una clasificación sostenible y comprensible. \evid{S07}{7}
\item \textbf{C9:} Registraron seis propuestas y dos modificaciones, con preguntas y aportes sobre inventario virtual y mantenimiento del orden. \evid{S07}{8}
\item \textbf{C10:} Contrastaron señalética e inventario web, con disposición estudiantil pendiente y una celda superpuesta cuya lectura se limita al cierre narrativo. \evid{S07}{9}
\end{compactitems} & La mesa de objetos perdidos se desorganiza y dificulta encontrar pertenencias. & \begin{compactitems}
\item Proponen un inventario virtual en la página del colegio. \evid{S07}{8}
\item Proponen estudiantes responsables por rotación para ayudar a organizar la mesa. \evid{S07}{8}
\end{compactitems} \\
\midrule
6 & \begin{compactitems}
\item Sebastián Lozano Barreneche
\item Alejandro Camargo Roldan
\end{compactitems} & \begin{compactitems}
\item \textbf{C5:} Integraron el recorrido y cinco hallazgos sobre espera, descanso, colados, momento de compra y acceso a dinero o tiquetes. \evid{S08}{9--10}
\item \textbf{C6:} Priorizaron provisionalmente la cafetería con 7 puntos; las alternativas sobre agua y baños quedaron sin datos. \evid{S08}{11--12}
\item \textbf{C7:} Reformularon el desafío para comprar durante el descanso sin perderlo ni llegar tarde, sin atribuir toda la demora a la fila. \evid{S08}{13--14}
\item \textbf{C8:} Conectaron decisiones económicas bajo incertidumbre y universalización de Kant para examinar compra y respeto del turno. \evid{S08}{14}
\item \textbf{C9:} Registraron seis ideas y dos combinaciones con preguntas cruzadas sobre autopedido, entrega separada y pedidos numerados. \evid{S08}{15--17}
\item \textbf{C10:} Contrastaron dos propuestas con diferencias de pago y compra, ajustaron carnet o tiquete y conservaron pedidos en fila. \evid{S08}{17--18}
\end{compactitems} & Filas y colados dificultan comer durante el descanso y llegar a clase a tiempo. & \begin{compactitems}
\item Proponen autopedido con carnet, tiquete o efectivo y recogida separada, aún por verificar. \evid{S08}{16, 18}
\item Proponen pedidos anticipados numerados, conservando la opción de pedir en fila. \evid{S08}{16, 18}
\end{compactitems} \\
\midrule
7 & \begin{compactitems}
\item Pablo José Del Cairo Bolaños
\item Nicolás Clavijo González
\end{compactitems} & \begin{compactitems}
\item \textbf{C5:} Relacionaron usuarios y recorrido de residuos con siete hallazgos y un supuesto, apoyados en observaciones y diez encuestas. \evid{S09}{4--10}
\item \textbf{C6:} Compararon clasificación, acceso a canecas y limpieza; desempataron dos opciones de 8 puntos a favor de aprender a clasificar por viabilidad. \evid{S09}{11--12}
\item \textbf{C7:} Delimitaron el desafío a identificar la caneca adecuada para estudiantes de sexto a undécimo en los descansos. \evid{S09}{12}
\item \textbf{C8:} Conectaron significantes y código de colores; quedó pendiente verificar rótulos y colores del patio. \evid{S09}{13}
\item \textbf{C9:} Registraron seis ideas, dos combinaciones y aportes cruzados sobre ubicación de canecas y continuidad del mensaje visual. \evid{S09}{14--15}
\item \textbf{C10:} Contrastaron información visual y agrupación de canecas, manteniendo señales rápidas y acceso cercano sin desatender zonas. \evid{S09}{15}
\end{compactitems} & Residuos mezclados, basura en el piso y dudas de clasificación aumentan el trabajo de aseo. & \begin{compactitems}
\item Proponen la misma imagen y color en empaque y caneca para guiar la clasificación. \evid{S09}{14}
\item Proponen guías estudiantiles rotativos por grado, con comentarios semanales del PAE. \evid{S09}{14}
\end{compactitems} \\
\midrule
8 & \begin{compactitems}
\item Andrés Camilo Maldonado Lemus
\item Carlos Samuel López Castrillón
\item Thomas Enrique Díaz Garzón
\end{compactitems} & \begin{compactitems}
\item \textbf{C5:} Organizaron usuarios, recorrido y cuatro hallazgos sobre residuos en graderías a partir de treinta encuestas. \evid{S10}{1--2}; \evid{S11}{1--2}
\item \textbf{C6:} Compararon recipientes, cultura ciudadana y empaques, y priorizaron acceso a recipientes con 9 puntos, limitado por mantenimiento. \evid{S10}{2--3}
\item \textbf{C7:} La verificación paginada está pendiente: S12 es un DOCX accesible, sin PDF de esta etapa entre las fuentes enlazadas. \srclink{S12}
\item \textbf{C8:} La verificación paginada está pendiente por el formato DOCX de S12; los PDF revisados solo cubren encuesta y C5--C6. \srclink{S12}
\item \textbf{C9:} La verificación paginada está pendiente: S13 es un DOCX accesible, lo que no implica ausencia de trabajo. \srclink{S13}
\item \textbf{C10:} La verificación paginada está pendiente porque S13 no aporta una versión PDF de esta etapa. \srclink{S13}
\end{compactitems} & Acumulación recurrente de envolturas y envases en las graderías, relacionada con hábitos y acceso a canecas. & \begin{compactitems}
\item No se identificó un conjunto de soluciones en los PDF; S10 aún no diseña solución y las posibles ideas de S13 siguen pendientes de verificación paginada. \evid{S10}{3}; \evid{S11}{1--2}; \srclink{S13}
\end{compactitems}
{\footnotesize C7--C8: \srclink{S12}.} \\
\midrule
9 & \begin{compactitems}
\item Sofía Guampe Granados
\item Nicolás Ospina Ramírez
\end{compactitems} & \begin{compactitems}
\item \textbf{C5:} Organizaron recorrido y cinco hallazgos con quince encuestas y observaciones, reconociendo fechas y mediciones faltantes. \evid{S14}{1--5}
\item \textbf{C6:} Compararon espera y orden; eligieron provisionalmente reducir la espera del almuerzo, con viabilidad sin datos. \evid{S14}{6--7}
\item \textbf{C7:} Revisaron el desafío para reducir la espera de bachillerato al comprar o reclamar almuerzo, sin imponer un medio específico. \evid{S14}{8}
\item \textbf{C8:} Relacionaron población y muestra con participación y escucha como apoyo; quedó pendiente conocer al personal. \evid{S14}{9}
\item \textbf{C9:} Presentaron seis ideas y dos combinaciones para explorar; preguntas y aportes figuran como propuestas, sin participación realizada verificable. \evid{S14}{10--11}
\item \textbf{C10:} Examinaron pedidos anticipados e información visible mediante supuestos, condiciones de acceso y preguntas sobre demora y funcionamiento. \evid{S14}{12--13}
\end{compactitems} & Las esperas dificultan comprar, comer y regresar a clase a tiempo. & \begin{compactitems}
\item Proponen explorar pedidos anticipados con reclamo diferenciado e información clara, sin exigir aplicación ni salida anticipada. \evid{S14}{10--13}
\item Proponen turnos por llegada, acuerdos claros y consulta a usuarios. \evid{S14}{11}
\end{compactitems} \\
\midrule
10 & \begin{compactitems}
\item Julián Alejandro Maldonado Baquero
\end{compactitems} & \begin{compactitems}
\item \textbf{C5:} Describió olvidos de materiales y entregas con perspectivas de estudiante y profesor sobre carga académica, prioridades y registro. \evid{S15}{1--2}
\item \textbf{C6:} Comparó priorizar y anotar tareas; eligió registrar pendientes por 8,5 a 7,5, condicionado a consultar y actualizar las anotaciones. \evid{S15}{2--4}
\item \textbf{C7:} Formuló un desafío sobre herramientas para evitar olvidos de bachillerato en espacios de estudio, ligado al registro de responsabilidades. \evid{S15}{5}
\item \textbf{C8:} Relacionó psicología cognitiva y UX/UI para comprender memoria y atención y orientar consultas sencillas de pendientes. \evid{S15}{5--6}
\item \textbf{C9:} Propuso mini tablero, chatbot y web, y atribuyó a ChatGPT la discusión y revisión de las dos propuestas mejoradas. \evid{S15}{6--7}
\item \textbf{C10:} Analizó tablero y chatbot-web con supuestos de consulta y pendientes sobre ubicación, actualización y acceso a WhatsApp e internet. \evid{S15}{7--8}
\end{compactitems} & Olvidos de trabajos, fechas y materiales escolares, con préstamos de compañeros como respuesta temporal. & \begin{compactitems}
\item Propone un mini tablero visual escolar con pendientes y retiro de actividades cumplidas, útil sin celular. \evid{S15}{7}
\item Propone chatbot de WhatsApp con recordatorios y enlaces a una web de tareas, materiales y fechas. \evid{S15}{7}
\end{compactitems} \\
\midrule
11 & \begin{compactitems}
\item Santiago Ospina Ramírez
\item Juan José Méndez Ramos
\end{compactitems} & \begin{compactitems}
\item \textbf{C5:} Organizaron cuatro hallazgos sobre hábitos, coste y oferta con diez observados en una fecha, sin generalización, acciones individuales ni citas textuales. \evid{S16}{1--5, 10--11}
\item \textbf{C6:} Compararon oferta en cafetería, barreras económicas y alimentación diaria, y priorizaron provisionalmente la primera con datos faltantes. \evid{S16}{5--11}
\item \textbf{C7:} Delimitaron el desafío al acceso a comida saludable y asequible en cafetería durante los descansos. \evid{S17}{1--3}
\item \textbf{C8:} Conectaron precios, accesibilidad y nutrición, señalando materiales consultados y datos de oferta, precios y preferencias faltantes. \evid{S17}{3--5}
\item \textbf{C9:} Registraron seis ideas y dos combinaciones con preguntas y aportes; algunas propuestas exceden el alcance del desafío. \evid{S18}{1--4}
\item \textbf{C10:} Contrastaron menú ampliado y alternativas menos procesadas con supuestos de compra, precios y recursos, sin aceptación demostrada. \evid{S18}{4--6}
\end{compactitems} & Dificultades para acceder a comida saludable y asequible en la cafetería durante los descansos. & \begin{compactitems}
\item Proponen ampliar el menú saludable, variado y asequible; aceptación no demostrada. \evid{S18}{3, 5--6}
\item Proponen alternativas menos procesadas y campañas sobre agua y verduras; preferencia específica no documentada. \evid{S18}{3, 5--6}
\end{compactitems} \\
\midrule
12 & \begin{compactitems}
\item Pedro Martín Ramírez Barrero
\item Santiago Tapias López
\end{compactitems} & \begin{compactitems}
\item \textbf{C5:} Organizaron perfiles, testimonios y una comparación de olores, limpieza, daños y necesidades de usuarios de los baños. \evid{S19}{1--3}
\item \textbf{C6:} Identificaron la necesidad de baños limpios y en buen estado; no se identificó matriz comparativa de oportunidades en el PDF. \evid{S19}{1--4}
\item \textbf{C7:} Formularon y revisaron el desafío para mejorar limpieza, olores y mantenimiento de baños usados por décimo durante la jornada. \evid{S19}{3}
\item \textbf{C8:} Relacionaron higiene, microorganismos y saneamiento con responsabilidad colectiva y cuidado de espacios comunes. \evid{S19}{3}
\item \textbf{C9:} No se identificó registro de ideación en el PDF legible, que finaliza con C7--C8. \evid{S19}{1--4}
\item \textbf{C10:} No se identificó contraste de ideas con evidencia del usuario en el PDF legible revisado. \evid{S19}{1--4}
\end{compactitems} & Malos olores, suciedad y daños generan incomodidad y evitan que algunos estudiantes usen los baños. & \begin{compactitems}
\item Plantean mejorar control y horarios de aseo como posibilidad, sin implementación documentada. \evid{S19}{2--3}
\item Plantean reparar sanitarios, lavamanos y otras instalaciones deterioradas. \evid{S19}{2}
\end{compactitems} \\
\midrule
13 & \begin{compactitems}
\item Tomás Mauricio Rincón Castañeda
\end{compactitems} & \begin{compactitems}
\item \textbf{C5:} Integró el recorrido de compra y reclamo con tres hallazgos sobre colados, saturación del personal y pérdida del descanso. \evid{S20}{1--3}
\item \textbf{C6:} Comparó tres oportunidades e integró orden y prevención de colados con delimitación y acompañamiento, sujetos a permisos y disponibilidad. \evid{S20}{3--5}
\item \textbf{C7:} Retiró barandas y vigilancia de la pregunta y delimitó el desafío a evitar colados y ordenar compra y reclamo en ambos descansos. \evid{S20}{5}
\item \textbf{C8:} Relacionó ciudadanía y autorregulación con escasez, demanda, coste de oportunidad e incentivos, separando conocimientos de hipótesis de aceptación. \evid{S20}{6}
\item \textbf{C9:} Registró tres ideas y dos ampliaciones apoyadas en ciudadanía y economía. \evid{S20}{7--8}
\item \textbf{C10:} Contrastó supervisión y fila rápida, ajustó turnos y pago físico, y dejó pendientes disposición de monitores y capacidad para separar flujos. \evid{S20}{8--10}
\end{compactitems} & Colados y congestión en compra y reclamo saturan al personal y consumen el descanso. & \begin{compactitems}
\item Propone canal único demarcado con cintas móviles y monitor en la entrada. \evid{S20}{8}
\item Propone fila rápida para preempacados con efectivo exacto o ficha física, separada de almuerzos. \evid{S20}{8--10}
\end{compactitems} \\
\end{longtable}
\end{document}
